\section{Opis pakietu danych badawczych}\label{sec:pakiet_danych}

Korpus źródłowy obejmuje 72 schematyzmy diecezjalne z 23 diecezji polskich, reprezentujące dwa główne przedziały chronologiczne: lata sześćdziesiąte-osiemdziesiąte XIX wieku oraz lata trzydzieste XX wieku.
Każdy schematyzm został zdigitalizowany w postaci pliku PDF, a następnie podzielony na osobne obrazy stron w wysokiej rozdzielczości, stanowiące podstawowe dane wejściowe systemu ekstrakcji.

Dane referencyjne, względem których oceniana jest jakość automatycznej ekstrakcji, stanowi baza danych zbudowana ręcznie przez zespół historyków w ramach projektu <tutaj dac nazwe projektu>.
Dla każdego schematyzmu baza zawiera sekwencję rekordów parafialnych z wypełnionymi polami: dekanat, parafia, wezwanie i materiał budowlany, przyporządkowanymi do konkretnych stron źródłowych.
Zbiór ten, liczący łącznie ok. 35, 000 rekordów, pełni podwójną rolę: jest zarówno docelowym produktem prac dokumentacyjnych projektu głównego, jak i wzorcem (\textit{ground truth}) do ewaluacji narzędzia cyfrowego.